\selectlanguage{ngerman}%
\frenchspacing%
\begin{ledgroupsized}[r]{120mm}%
\footnotesize%
\pstart%
\noindent\textbf{Überlieferung:}%
\pend%
\end{ledgroupsized}%
\begin{ledgroupsized}[r]{114mm}%
\footnotesize%
\pstart%
\parindent -6mm%
\makebox[6mm][l]{\textit{L}}%
Auszug aus
\textsc{E.\;Weigel} \bracks{Praeses}, \textit{Pendulum ex tetracty deductum}
\bracks{Dissertation: Kandidat \textsc{H.\:C. Ebell}},
Jena 1674 (nicht paginiert):\cite{02541}
LBr~250 Bl.\;1.
Ein Papierstreifen.
Drei Zeilen auf Bl.\;\href{\Scan{LBr_250_001r}}{1\,r\textsuperscript{o}};
die \href{\Scan{LBr_250_001v}}{Rückseite} ist unbeschrieben.
\pend%
\end{ledgroupsized}%
\vspace{5mm}%
\begin{ledgroup}%
\footnotesize%
\pstart%
\noindent%
\textbf{Datierungsgründe:}
Der Träger des undatierten, notizartigen Auszugs N.\;\ref{38575} weist weder Spuren eines Wasserzeichens noch irgendein äußeres datierungsrelevantes Merkmal auf.
Die im Text zitierte Dissertation \textit{Pendulum ex tetracty deductum}\cite{02541} wurde im Herbst 1674 in Jena\indexO{Jena} veröffentlicht (das Vorwort des Betreuers Erhard Weigel\indexN{\textls{Weigel}, Erhard 1625\textendash1699} trägt das Datum 21.\;September 1674).
Da sich Leibniz zu jener Zeit in Paris\indexO{Paris} aufhielt, ist anzunehmen, dass er die Dissertation erst nach seiner Übersiedlung nach Hannover\indexO{Hannover} Mitte Dezember 1676 lesen konnte.
Hieraus ergibt sich der Terminus post quem der vorgeschlagenen Datierung.%
\pend%
\pstart%
Relevant für die Bestimmung des Terminus ante quem ist möglicherweise die Tatsache, dass der Auszug N.\;\ref{38575} -- lediglich eine wohlgesinnte Anspielung auf Leibnizens \textit{Hypothesis physica nova}\cite{00256} umfassend -- im Nachlass unter der Signatur LBr~250 überliefert ist, welche den Briefwechsel zwischen Leibniz und Honoré Fabri\indexN{\textls{Fabri}, Honoré 1607\textendash1688} umfasst.
Leibnizens letzter Brief an Fabri, wahrscheinlich in den ersten Monaten 1677 verfasst (vgl. \textit{LSB} II,\;1\textsuperscript{2} N.\;133),\cite{01309} befasst sich unter anderem mit der aus der \textit{Hypothesis physica nova}\cite{00256} stammenden Erklärung der Natur physikalischer Körper, auf welche auch der Auszug N.\;\ref{38575} Bezug nimmt.
Diese Feststellung lässt die Annahme zu, dass N.\;\ref{38575} im Zusammenhang mit Leibnizens letztem Brief an Fabri\indexN{\textls{Fabri}, Honoré 1607\textendash1688} entstanden sein könnte.
Eine spätere Datierung ist jedoch möglich, zumal Leibniz sich noch lange mit dem Erklärungsmodell aus der \textit{Hypothesis physica nova}\cite{00256} befasst hat, wie etwa einige in \textit{LSB} VIII,\;3 veröffentlichte Texte zeigen.%
\pend%
\end{ledgroup}%
\selectlanguage{latin}%
\frenchspacing%
\normalsize%
\vspace{8mm}%
\pstart%
\noindent%
\bracks{%
\href{\Scan{LBr_250_001r}}%
{1\,r\textsuperscript{o}}%
}
\edtext{%
Weigelius%
\indexN{\textls{Weigel}, Erhard 1625\textendash1699}
diss. de \textit{Pendulo}%
\indexS{pendulum}%
\indexS{dissertatio}
corollar. ultimo%
\bracks{:}
\textit{%
Speciosa est
\edtext{%
Hypothesis Leibnitiana,%
\indexS{hypothesis Leibnitiana}
quae bullulis%
\indexS{bullulae}
pleraque corporum phaenomena%
\indexS{phenomena corporum}
salvare docet.%
}{%
\lemma{\textit{Speciosa} \bracks{...} \textit{docet}}%
\Cfootnote{%
Anspielung an
\textsc{G.\:W.\;Leibniz},
\textit{Hypothesis physica nova},
§\;11\,ff.
(Mainz 1671;
\textit{LSB} VI,\;2 N.\;40, S.\;226\,ff.).%
\cite{00256}%
}}%
}%
}{%
\lemma{Weigelius \bracks{...} \textit{docet}}%
\Cfootnote{%
\textsc{E.\;Weigel} \bracks{Präses}, \textit{Pendulum ex tetracty deductum}
\bracks{Dissertation: Kandidat \textsc{H.\:C.\,Ebell}},
corollaria, n.\;8
(Jena 1674, Bl.\;\bracks{E4\,v\textsuperscript{o}}).
Die Vorlage liest \textit{Leibnüziana} statt \textit{Leibnitiana}.%
\cite{02541}%
}}
\pend%
\selectlanguage{latin}%
\frenchspacing%
\normalsize%
